% x: billion iterations per second, 0.37 cm per unit.
\begin{figure}[t]
\centering
\begin{tikzpicture}[font=\scriptsize,x=0.37cm,y=1cm,
  meas/.style={circle,fill=cA,inner sep=2.1pt},
  old/.style={circle,draw=cA,fill=white,thick,inner sep=1.8pt},
  ceil/.style={draw=cB,line width=1.6pt},
  band/.style={fill=cB!30,draw=cB},
  lab/.style={anchor=east,align=right,font=\scriptsize}]
% grid
\foreach \v in {0,5,10,15,20,25,30}{
  \draw[gridc] (\v,0.4) -- (\v,-6.2);
  \node[below,font=\scriptsize] at (\v,-6.2) {\v};
}
\node[below=10pt,font=\small] at (15,-6.2) {billion iterations per second, one device};
% Row 1: RTX PRO 6000 sigma
\node[lab] at (-0.4,0) {RTX PRO 6000\\$\sigma$-walk};
\draw[inkS!40,line width=3pt] (0,0) -- (14.64,0);
\node[meas] at (14.64,0) {}; \node[meas] at (15.88,0) {};
\node[above,font=\tiny] at (14.0,0.08) {14.64};
\node[above,font=\tiny] at (16.5,0.08) {15.88};
\draw[ceil] (23.97,-0.22) -- (23.97,0.22); \node[above,font=\tiny,align=center] at (23.97,0.22) {24.0 carry-less\\ceiling, rest free};
\fill[band] (24.7,-0.15) rectangle (26.1,0.15);
\node[right,font=\tiny,align=left] at (26.2,0) {24.7--26.1 one-add\\floor};
% Row 2: RTX PRO 6000 table
\node[lab] at (-0.4,-1.1) {RTX PRO 6000\\table walk};
\draw[inkS!40,line width=3pt] (0,-1.1) -- (8.24,-1.1);
\node[meas] at (8.24,-1.1) {}; \node[above,font=\tiny] at (8.24,-1.02) {8.24 exact (v3)};
\node[old] at (20.08,-1.1) {}; \node[above,font=\tiny] at (19.3,-1.02) {20.08 (v2)};
\draw[ceil] (22.0,-1.32) -- (22.0,-0.88); \node[below,font=\tiny] at (22.0,-1.32) {22.0 logic pipe};
\fill[band] (22.7,-1.25) rectangle (27.4,-0.95); \node[right,font=\tiny] at (27.5,-1.1) {22.7--27.4 SOL};
% Row 3: B200
\node[lab] at (-0.4,-2.2) {B200};
\draw[inkS!40,line width=3pt] (0,-2.2) -- (8.84,-2.2);
\node[meas] at (8.84,-2.2) {}; \node[above,font=\tiny] at (8.84,-2.12) {8.84 $\sigma$};
\node[old] at (19.40,-2.2) {}; \node[above,font=\tiny] at (19.40,-2.12) {19.40 table};
\draw[-{Stealth[length=4pt]},dashed,cB,thick] (19.9,-2.2) -- (29.6,-2.2);
\node[below,font=\tiny,align=center] at (24.8,-2.27) {30 needs a $1.55\times$ logic cut\\(carry-less floor: 255)};
% Row 4: RTX PRO 4500
\node[lab] at (-0.4,-3.3) {RTX PRO 4500\\(165\,W)};
\draw[inkS!40,line width=3pt] (0,-3.3) -- (5.107,-3.3);
\node[meas] at (5.107,-3.3) {}; \node[above,font=\tiny] at (4.6,-3.22) {5.11 $\sigma$};
\node[meas] at (12.50,-3.3) {}; \node[above,font=\tiny] at (12.5,-3.22) {12.50 bridge map};
\draw[ceil] (14.14,-3.52) -- (14.14,-3.08); \node[right,font=\tiny] at (14.3,-3.3) {14.1 map's arithmetic ceiling};
% Row 5: FPGA
\node[lab] at (-0.4,-4.4) {VU47P FPGA\\(AWS F2)};
\draw[inkS!40,line width=3pt] (0,-4.4) -- (8.25,-4.4);
\node[meas] at (8.25,-4.4) {}; \node[above,font=\tiny] at (8.0,-4.32) {8.25};
\fill[band] (8.8,-4.55) rectangle (9.3,-4.25); \node[right,font=\tiny] at (9.45,-4.4) {8.8--9.3 at the block-RAM limit (not met)};
% Row 6: NIST
\node[lab] at (-0.4,-5.5) {P-256 rho on\\RTX PRO 6000};
\draw[inkS!40,line width=3pt] (0,-5.5) -- (1.632,-5.5);
\node[meas] at (1.632,-5.5) {}; \node[above,font=\tiny] at (1.632,-5.42) {1.63};
\draw[ceil] (4.6,-5.72) -- (4.6,-5.28); \node[below,font=\tiny] at (4.6,-5.72) {4.6 multiplies only, $W{=}32$};
\draw[ceil] (12.2,-5.72) -- (12.2,-5.28); \node[right,font=\tiny] at (12.35,-5.5) {12.2 multiplies only, $W\to\infty$};
% key
\begin{scope}[shift={(0,-7.6)}]
\node[meas] at (0.3,0) {}; \node[anchor=west] at (0.8,0) {measured};
\node[old] at (6.3,0) {}; \node[anchor=west] at (6.8,0) {measured, superseded cycle rule};
\draw[ceil] (20.0,-0.2) -- (20.0,0.2); \node[anchor=west] at (20.4,0) {ceiling};
\fill[band] (24.0,-0.15) rectangle (25.4,0.15); \node[anchor=west] at (25.6,0) {band or projection};
\end{scope}
\end{tikzpicture}
\caption{How far each platform can go, against what it does today.
  Ceilings are derived from measured pipe rates with every other cost
  set to zero, so they are upper bounds, not targets: the RTX~PRO~6000's
  are its carry-less and logic pipes for any walk with one affine
  addition per step (Section~\ref{sec:floor}); the NIST bounds count only
  field multiplications at the measured multiply rate
  (Section~\ref{sec:bounds-nist}). None of these moves the $2^{60.9}$
  iteration count.}
\label{fig:bounds}
\end{figure}
